\section{固定电流维数以后，怎样评价一次材料更新}
\subsection{把要比较的问题固定下来}
电流空间的设计（current-space design）是一个模型选择问题。设接收几何和材料猜测已经确定，我们只允许切向电流方程保留 $k$ 个复电流方向。不同方法选择不同方向，然后各自计算材料步。公平比较要求材料状态、总场、照明、观测残差、物理材料坐标、正则化以及先验线性项（linear prior term）保持一致。这里没有给某个方法更好的初值，也没有把其材料步交给另一个新求解器修复。

令 $U$ 的列为物理电流内积下的正交方向，且压缩电流矩阵 $U^*LU$ 可逆；公式写在已经归一化的电流坐标中。材料正则化保证材料正规矩阵正定，却不保证这个电流矩阵可逆；实验按锁定阈值拒绝不稳定空间。$[U,C]$ 表示在原空间之外增加 $C$ 后重新正交化的增广空间，不是两个矩阵相加。Galerkin 电流模型（Galerkin current model）与材料导数为
\[
 R_U=U(U^*LU)^{-1}U^*,\qquad J_U=SR_UB.
\]
从右往左阅读：$B$ 把材料变化转为电流激励，$R_U$ 只在保留电流空间里解内部传播，$S$ 把结果变成接收场。$k$ 是所有照明（illumination）共享的实际复电流列数；它不是材料变量数，也不是给每次照明各分配一个独立的预算。

材料步 $s_U$ 最小化
\[
 \Phi_U(s)=\tfrac12\|r+J_Us\|^2+\ell^\top s+
 \tfrac12s^\top\Lambda s.
\]
$\ell$ 是当前先验线性项，$\Lambda$ 包含相同的正则化和阻尼；两者都来自共同状态。材料坐标 $s$ 为实数，包含复介电参数的实部和虚部；场仍为复数。复输出映射的材料伴随是实拉回（real pullback），即对复共轭转置的结果取实部。这一规则贯穿所有正规方程，不能把材料更新误当成任意复向量。

\subsection{为什么使用完整二次目标上的差距}
完整模型给出 $J=SL^{-1}B$、$H=J^*J+\Lambda$ 和完整最优步 $s_*$。约化模型在自己的目标上最优，不意味着它在完整目标上也好。因此把每个 $s_U$ 放回共同的完整二次目标 $\Phi$，定义完整二次差距（full-quadratic gap）
\[
 R(U)=\Phi(s_U)-\Phi(s_*) .
\]
正则化保证 $H\succ0$。在 $s_*$ 处梯度为零，把二次函数展开便得到
\[
 R(U)=\tfrac12(s_U-s_*)^\top H(s_U-s_*)
 =\tfrac12\|g(s_U)\|_{H^{-1}}^2,
\qquad g(s)=J^*(r+Js)+\Lambda s+\ell.
\]
这是同一目标中的恒等式：一步材料误差和完整目标的驻点缺陷是同一件事的两种测量方式。$H$ 能量考虑各材料方向的曲率，而体积加权材料范数衡量物理材料变化；两者用途不同。这个指标严格说明局部更新的质量，还没有说明最后的非线性图像是否更准确。

\subsection{材料步差怎样从电流改变中产生}
设完整导数与当前导数相差 $D=J-J_U$，并记 $e_U=r+J_Us_U$。约化步满足
\[
 J_U^*e_U+\Lambda s_U+\ell=0.
\]
代入完整梯度，并利用 $J=J_U+D$，得到
\[
 \begin{aligned}
 g(s_U)&=(J_U+D)^*(e_U+Ds_U)+\Lambda s_U+\ell\\
 &=J_U^*Ds_U+D^*e_U+D^*Ds_U.
 \end{aligned}
\]
第一项描述遗漏响应与原有材料敏感度的耦合，第二项把原残差经遗漏导数拉回材料，第三项是遗漏导数本身的二次作用。三个项可能相互抵消或增强。导数差的范数没有给出它们与当前残差之间的方向和相位关系。

对于指定增广 $C$，用真实 Schur 因子
\[
 D_C=J_{[U,C]}-J_U=Y_C\Gamma_C^{-1}N_C
\]
替换 $D$，上述三项成为增广模型在旧步处的梯度 $h_C$。两个模型的正规方程相减，得
\[
 d_C=s_{[U,C]}-s_U=-H_{[U,C]}^{-1}h_C.
\]
这里完整保留 $\Gamma_C$ 的候选耦合。多照明先共享同一个材料向量，再把输出实化堆叠；一个复电流列也不因此必然对应一个秩一的实材料导数。

\subsection{有符号收益：改善方向与走得太远是两回事}
把 $d_C$ 代入同一个完整二次函数，精确有限收益（finite signed gain）为
\[
 Q_C=\Phi(s_U)-\Phi(s_U+d_C)
 =-g(s_U)^\top d_C-\tfrac12d_C^\top Hd_C.
\]
正值表示完整目标改善，负值表示变差。线性项检查是否朝向完整目标的下降方向；二次项计算这一步的曲率代价。忽略后者就是一阶有符号收益（first-order signed gain），只在其误差受到控制时才足以代替完整收益。

这不是新的通用优化恒等式。本文的电磁内容在于：$d_C$ 通过材料注入（material injection）、内部反馈（retained feedback）和接收返回（observation return）产生，且只能由实际可实现的电流空间生成。一般低秩材料估计没有这一实现约束。

\subsection{知道精确目标，与用有限信息估计它}
完整参考权限用于离线评价和最佳已找到搜索（best-found search），不提供给信息受限的选择器。实际有符号选择使用一个独立、更丰富的约化参考模型（richer ROM anchor）；有限探针版还支付少量完整切向及伴随作用（full tangent/adjoint actions）。如果参考模型的任务梯度不准确，精确的二次公式也不能保证其预测的排序准确。

通用正向—伴随电流降阶模型（generic primal--dual current ROM）通过本征正交分解（proper orthogonal decomposition，POD）压缩相应电流样本，也是 task-oriented reduction 的一种适配。它从同一信息预算构造实际复电流 POD 空间。与它比较，是为了判断反馈和有限步评分是否提供额外价值，而不是把已有的伴随思想换个名字。

各方法必须达到规定的实际 $k$；若达不到，记录不可行（infeasible）。因此逐次选择中可能出现负分数，但还需加入方向才能达到固定预算。负分数不是“算法保证下降”的证据，也不独立证明最终空间一定差。完整结果需要在同一个完整二次目标上检验，并保留所有对象和预算。
